\documentclass[11pt,a4paper]{article}
\usepackage[T1]{fontenc}
\usepackage[utf8]{inputenc}
\usepackage[french]{babel}
\usepackage{amsmath,amssymb,booktabs,longtable,array}
\usepackage[a4paper,margin=2.4cm]{geometry}
\usepackage{xcolor}
\usepackage[hidelinks]{hyperref}
\usepackage{microtype}
\setlength{\parindent}{0pt}
\setlength{\parskip}{0.55em}
\newcommand{\NW}{\mathrm{NW}}
\newcommand{\dd}{\mathrm{d}}
\newcommand{\hyp}{\textcolor{orange!70!black}{Hypothèse}}
\newcommand{\plaus}{\textcolor{blue!65!black}{Plausible}}
\newcommand{\obs}{\textcolor{green!45!black}{Observé}}
\newcommand{\und}{\textcolor{red!65!black}{Non démontré}}

\title{Lecture critique du modèle M2\\
\large Distribution Boltzmann--Pareto et auto-organisation critique}
\author{Réplication computationnelle indépendante}
\date{3 juillet 2026}

\begin{document}
\maketitle

\begin{abstract}
Le document de conception M2 propose un modèle parcimonieux de capital réel,
crédit nominal perpétuel, chocs multiplicatifs, faillites en cascade et
population ouverte. Sa contribution la plus solide est de transformer plusieurs
mécanismes opaques du WIP en règles testables et d'identifier explicitement
l'effet cohorte comme menace. En revanche, ni le corps exponentiel ni
l'auto-organisation de l'exposant de Pareto ne sont démontrés. Les explorations
E0--E7 constituent des ablations informatives mais non une validation
indépendante. Plusieurs règles de faillite et de mesure restent sous-spécifiées.
La conclusion épistémique est donc : modèle cohérent comme programme de
recherche, mécanisme Boltzmann conjectural, queue de Pareto préliminaire, SOC non
établie.
\end{abstract}

\tableofcontents

\section{Objectifs, méthode et statut des preuves}

Cette lecture reconstruit fidèlement M2, sépare faits, arguments et conjectures,
puis transforme les ambiguïtés en exigences d'implémentation ou expériences de
réfutation. La source principale est
\texttt{recherche/conception\_boltzmann\_pareto\_soc/conception\_modele\_M2.md},
lue intégralement avec ses annexes. Le prototype, les sondes et les logs E0--E7
sont des éléments produits dans le même cycle de conception : ils soutiennent
la plausibilité, mais ne sont ni une réplication ni une preuve analytique.

\begin{longtable}{p{0.18\textwidth}p{0.31\textwidth}p{0.43\textwidth}}
\toprule
Statut & Contenu & Évaluation \\
\midrule
\endhead
\obs & Conservation du principal à la création d'un prêt ; espérance unitaire du
choc log-normal ; point fixe autarcique. & Conséquences algébriques des règles,
sous réserve d'une implémentation correcte. \\
\plaus & Population bornée dans E7c ; queue compatible avec Pareto ; faible
corrélation âge--richesse. & Observations préliminaires sur une configuration et
un petit nombre de traces ; à répliquer par seed et fenêtre. \\
\hyp & Flux du corps indépendants de $\NW$ ; stabilisation de l'exposant par les
cascades ; extensivité en $\lambda$. & Hypothèses mécanistiques falsifiables. \\
\und & Corps exponentiel, SOC au sens critique, absence de réglage fin,
stationnarité asymptotique. & Aucun théorème ni campagne suffisante dans le
rapport. \\
\bottomrule
\end{longtable}

\section{Résumé fidèle du modèle M2}

\subsection{État et observables}

Une entité vivante $i$ porte un capital réel scalaire $w_i\geq0$, une date de
naissance et une dette d'amorçage nominale universelle $d_0$. Un contrat actif
$e=(\ell,b,q,r)$ relie un prêteur $\ell$ à un emprunteur $b$, avec principal
nominal constant $q>0$ et taux par pas $r\geq0$. Les prêts sont perpétuels : il
n'existe ni échéance ni amortissement.

En notant
\[
 C_i=\sum_{e:\ell(e)=i}q_e,\qquad
 D_i=\sum_{e:b(e)=i}q_e,
\]
la variable cible est
\[
 \NW_i=w_i+C_i-D_i-d_0.
\]
Le capital brut $w_i$ est diagnostique. Le rapport définit le revenu brut par
$y_i=\alpha\sqrt{w_i}+I_i^{\text{reçu}}$ et le revenu net en retranchant
$I_i^{\text{payé}}$. Cette définition n'inclut ni plus-value de $w$, ni pertes sur
créances, ni variation de valeur nette ; elle devra donc être nommée
\emph{revenu de flux} pour éviter une équivalence abusive avec $\Delta\NW$.

\subsection{Séquence complète d'un pas}

L'ordre proposé est constitutif du modèle :
\begin{enumerate}
  \item naissances Poisson, chaque entrante recevant $(w_0,d_0)$ et
  $\NW_0=\varepsilon_0=w_0-d_0>0$ ; $N(0)=0$ ;
  \item choc multiplicatif indépendant sur chaque capital réel ;
  \item extraction concave ajoutée au capital ;
  \item service des intérêts contractuels et marquage des défauts ;
  \item dépréciation du seul capital réel ;
  \item $\lfloor N/k\rfloor$ rencontres de crédit locales ;
  \item résolution itérative des défauts et insolvabilités jusqu'à stabilité.
\end{enumerate}

\begin{figure}[ht]
\centering
\fbox{\begin{minipage}{0.92\textwidth}\centering
Naissances $\rightarrow$ choc sur $w$ $\rightarrow$ extraction
$\rightarrow$ intérêts $\rightarrow$ dépréciation $\rightarrow$ crédit
$\rightarrow$ cascade de faillites $\rightarrow$ mesure
\end{minipage}}
\caption{Ordre des événements de M2. Le modifier définit une autre dynamique
discrète, surtout pour les nouveaux prêts et les insolvabilités temporaires.}
\end{figure}

\subsection{Population ouverte et paramètres}

Les naissances suivent $n_t\sim\mathrm{Poisson}(\lambda)$ et la mort résulte
uniquement du défaut ou de $\NW<0$. Il n'existe ni capacité de charge ni probabilité
de décès exogène. Le bornage allégué est $N^*=\lambda\,\mathbb{E}[L]$ si la durée
de vie $L$ a effectivement une espérance finie pour toutes les positions du
réseau. C'est une conséquence conditionnelle, pas la preuve que cette espérance
est finie.

Les valeurs dynamiques sont $\delta,\lambda,k,\sigma,w_0,d_0$ avec $\alpha=1$
comme convention. Dire « six paramètres » masque que $(w_0,d_0)$ représentent
deux degrés de liberté, souvent reformulés en $(w_0,\varepsilon_0/w_0)$.
$\delta$ et $\lambda$ ne sont des conventions d'échelle qu'après démonstration
des invariances de temps et d'extensivité.

\section{Reconstruction mathématique}

\subsection{Dynamique du capital hors crédit}

Pour chaque entité vivante, le choc est
\[
 \eta_{i,t}\sim\mathcal N(-\sigma^2/2,\sigma^2),\qquad
 w_i\leftarrow w_i\exp(\eta_{i,t}),
 \quad\mathbb E[\exp(\eta)]=1.
\]
Puis l'extraction et la dépréciation donnent successivement
\[
 w_i\leftarrow w_i+\alpha\sqrt{w_i},\qquad
 w_i\leftarrow(1-\delta)w_i.
\]
Sans choc ni crédit, la carte exacte d'un pas n'a pas exactement le même point
fixe que l'ODE $\dot w=\alpha\sqrt w-\delta w$, car extraction et dépréciation
sont séquentielles. L'ODE a $w_+=(\alpha/\delta)^2$, tandis que la carte écrite
ci-dessus satisfait
\[
 w=(1-\delta)(w+\alpha\sqrt w)
 \quad\Longrightarrow\quad
 w_+=\left(\frac{(1-\delta)\alpha}{\delta}\right)^2.
\]
Pour $\delta=0{,}05$, l'écart est d'environ 9,75\,\% en niveau. Employer 400
comme point fixe exact est donc une approximation de temps continu.

\subsection{Marché du crédit}

Le rendement marginal de l'extraction est
\[
 r_i^*=\frac{\dd(\alpha\sqrt w)}{\dd w}\bigg|_{w_i}
      =\frac{\alpha}{2\sqrt{w_i}}.
\]
Dans un échantillon uniforme de $k$ vivantes, $\ell$ maximise $w$ et $b$ le
minimise. Pour $w_\ell,w_b>0$,
\[
 r=\sqrt{r_\ell^*r_b^*},\qquad \rho=r+\delta,
 \qquad w^*(\rho)=\left(\frac{\alpha}{2\rho}\right)^2.
\]
L'offre, la demande et le volume sont
\[
 q_{\rm dem}=[w^*-w_b]_+,
 \quad q_{\rm off}=[w_\ell-w^*]_+,
 \quad q=\min(q_{\rm dem},q_{\rm off}).
\]
Si $q>0$, le transfert est conservatif :
\[
 w_\ell' = w_\ell-q,\qquad w_b'=w_b+q,
 \qquad w_\ell'+w_b'=w_\ell+w_b,
\]
et un contrat $(\ell,b,q,r)$ est créé. La valeur nette de chacun est également
inchangée à la création : le prêteur remplace $q$ de réel par une créance, et
l'emprunteur reçoit $q$ de réel avec une dette égale.

Le cas $w=0$ rend $r_i^*$ infini et n'est pas défini dans le rapport. Une limite
analytique conduit généralement à $w^*=0$ et $q=0$, mais une régularisation
numérique choisie autrement introduirait un paramètre caché.

\subsection{Intérêts, défaut et conservation}

Pour un contrat $e$, l'échéance du pas est $a_e=r_eq_e$. Un paiement isolé peut
s'écrire
\[
 p_e=\min(w_b,a_e),\qquad
 w_b\leftarrow w_b-p_e,\quad w_\ell\leftarrow w_\ell+p_e.
\]
Ainsi le capital total est conservé même lors d'un paiement partiel ; c'est la
créance d'intérêt impayée qui n'est pas capitalisée. Si $p_e<a_e$, $b$ est en
défaut et doit être liquidé en phase 7. Traiter les contrats séquentiellement
donne une priorité au premier contrat rencontré. Sans ordre canonique ou service
proportionnel simultané, la permutation du carnet modifie les gagnants et les
cascades.

\subsection{Faillite et cascade}

Soit $i$ en défaut ou avec $\NW_i<0$. Le texte exige : annuler ses dettes
contractuelles ; transférer les créances qu'elle détient et son capital résiduel
à ses créanciers au prorata ; éteindre $d_0$ ; retirer $i$ ; puis recommencer
jusqu'à ce qu'aucune vivante ne vérifie le critère.

Une formalisation naturelle, mais non entièrement spécifiée, pose pour chaque
créancier $j$ de $i$
\[
 s_{ji}=\frac{q_{ji}}{D_i},\qquad
 w_j\leftarrow w_j+s_{ji}w_i,
\]
et remplace chaque créance $(i,h,q_{ih},r_{ih})$ par des fractions
$(j,h,s_{ji}q_{ih},r_{ih})$. Cette règle exige de fractionner les contrats. Elle
ne définit pas $D_i=0$, les créanciers déjà promis à la faillite, les
auto-contrats créés par transfert, ni l'ordre entre plusieurs faillites
simultanées. Ces cas peuvent changer la contagion et doivent être spécifiés avant
codage.

\section{Hypothèses fortes et paramètres cachés}

\begin{longtable}{p{0.31\textwidth}p{0.27\textwidth}p{0.34\textwidth}}
\toprule
Proposition & Statut & Test discriminant \\
\midrule
\endhead
Le corps de $\NW$ est exponentiel. & \und ; le rapport signale que gamma ou
log-normale gagnent encore. & AIC/BIC sur domaine déclaré, QQ-plot, stabilité aux
seuils de corps, $T$ contre $B_0/A_0$. \\
L'exposant de Pareto est stabilisé par une boucle SOC. & \hyp ; stabilité
temporelle ne prouve pas criticité. & Multi-seed, $x_{\min}$ libre et commun,
ablation du crédit/cascades, finite-size scaling, réponse à $k$ et $\sigma$. \\
Les cascades suffisent à renouveler les classes. & \plaus dans E7c seulement. &
Mortalité et entrées du top décile, matrices de transition par quantile, risque
conditionnel à âge et richesse. \\
Les classes ne sont pas des cohortes. & \plaus, non acquis. & Distribution
jointe âge--$\NW$, queues dans tranches d'âge, top renewal, réplication après
transitoire long. \\
La population est bornée endogènement. & \plaus si toute durée de vie a une
espérance finie. & Pente de $N(t)$, équilibre naissances--décès, durée de vie
censurée, survie du top, horizon doublé. \\
La forme est robuste sans réglage fin. & \und ; E7 passe par une « fenêtre de
naissance ». & Grille pré-enregistrée incluant les échecs, plusieurs seeds, cartes
continues autour des frontières. \\
$\NW$ est la variable cible économiquement pertinente. & Choix de modélisation.
& Comparer $w$, $\NW$, revenu brut/net et $\Delta\NW$ ; tester la sensibilité à
$d_0$. \\
La moyenne géométrique du taux est neutre. & Convention arbitraire. & Ablations
géométrique, arithmétique et négociation bornée, sans recalibrage. \\
Le transfert prorata est neutre. & Convention arbitraire et potentiellement
redistributive vers les gros créanciers. & Annulation, destruction, transfert
proportionnel et enchère ; comparer concentration et cascades. \\
\bottomrule
\end{longtable}

Les paramètres cachés comprennent aussi le pas de temps, l'ordre des phases,
l'échantillonnage avec ou sans remise dans un round, la possibilité de participer
plusieurs fois par pas, l'ordre du carnet lors du service, la tolérance de
solvabilité, le nombre minimal de points de queue, le quantile de corps (95\,\%),
la cadence des snapshots, le nombre de réplications bootstrap et le traitement
des décès simultanés. Ce sont des choix structurels ou statistiques, pas de
simples détails logiciels.

\section{Critique scientifique structurée}

\subsection{Le modèle est-il vraiment minimal ?}

Il est nettement plus minimal que le WIP comptable : un stock réel, un réseau de
contrats et un seuil de solvabilité. Il n'est toutefois pas minimal au sens
causal. Trois mécanismes produisent déjà de la multiplicativité (GBM, intérêts
proportionnels, pertes proportionnelles) et deux mécanismes imposent un bord
(faillite, $d_0$). Sans ablations factorielles, on ne sait pas lequel est
nécessaire à la queue. De plus, une dette $d_0$ sans créancier ni remboursement
est économiquement un seuil de liquidation individuel, pas une dette au sens
comptable. La nommer dette donne une microfondation que la règle ne possède pas.

\subsection{Règles contestables}

Les prêts perpétuels sans remboursement permettent des paiements d'intérêts
cumulés arbitrairement supérieurs au principal. La cible $w^*(r+\delta)$ suppose
qu'un taux contractuel et la dépréciation du capital s'additionnent comme coûts
marginalement comparables, alors que l'emprunteur garde le capital et ne rembourse
jamais le principal. Le prêteur compare par ailleurs le rendement marginal perdu
à un taux sans intégrer explicitement risque de défaut ou valeur de récupération.
La règle max/min parmi $k$ sélectionne mécaniquement des écarts extrêmes et peut
fabriquer une stratification. Enfin, transférer les actifs de la faillie aux
créanciers concentre potentiellement le réseau au lieu de le désenchevêtrer.

\subsection{Mécanismes susceptibles de fabriquer une queue de Pareto}

Le choc log-normal multiplicatif avec réinjection et absorption suffit souvent à
générer un processus de Kesten. Une queue observée ne prouve donc ni le rôle du
crédit ni la SOC. Les prêts perpétuels cumulent des créances proportionnelles à
la taille ; le mécanisme max/min et les transferts prorata favorisent aussi les
agents déjà riches. Enfin, optimiser $x_{\min}$ sur plusieurs seuils puis retenir
la meilleure lecture crée un biais de sélection. Les contrôles indispensables
sont GBM seul avec le même mécanisme démographique, crédit sans choc, faillites
sans transfert de créances et réseau rebrassé conservant les bilans.

\subsection{Mécanismes susceptibles d'empêcher le corps exponentiel}

L'extraction $\sqrt w$ et la dépréciation créent un rappel vers un niveau positif,
donc un mode intérieur plutôt qu'en zéro. Le choc sur $w$, alors que le corps est
analysé sur $\NW=w+C-D-d_0$, n'est pas additif conditionnellement à $\NW$ si les
compositions de bilan varient. Les entrants sont injectés en un point
$\varepsilon_0$, non par une distribution au bord. La population stationnaire est
un mélange d'âges et de positions de réseau. Le crédit n'est pas un échange
aléatoire de monnaie à somme conservée : il crée simultanément deux postes de
bilan et sélectionne les extrêmes. Aucune de ces propriétés ne garantit la
symétrie de renversement du temps invoquée pour Boltzmann.

\subsection{Les diagnostics proposés suffisent-ils ?}

Non. AIC/BIC et CSN sont nécessaires mais insuffisants sur des observations
temporellement dépendantes. Il faut des intervalles par blocs ou par trajectoire,
des tests prédictifs hors échantillon, des courbes de sensibilité au domaine
d'ajustement, le taux de faux positifs sur données synthétiques et des tailles
effectives d'échantillon. Pour une revendication SOC, une loi de puissance des
avalanches ne suffit pas : il faut au minimum vérifier absence de réglage du
point critique, séparation des échelles de forçage et relaxation, réponse à la
taille du système, et relations d'échelle robustes. Ici les naissances, chocs,
marché et cascades ont lieu au même pas, ce qui affaiblit l'analogie SOC.

\subsection{Tests supplémentaires requis}

Outre le protocole du rapport : permutations de l'ordre des contrats ; traitement
simultané des intérêts ; liquidation simultanée versus séquentielle ; test de
conservation global à chaque phase ; risque de décès conditionnel à $(\text{âge},
\NW,C/D)$ ; matrices de mobilité inter-déciles ; survie avec censure à droite ;
autocorrélation et taille effective des pools ; stabilité quand la fréquence de
snapshot change ; contrôle de taille en variant $\lambda$ ; ablations complètes
des trois sources multiplicatives ; et comparaison à des données synthétiques
log-normales connues.

\subsection{Où l'artefact de cohorte peut-il réapparaître ?}

$N(0)=0$ supprime une cohorte initiale explicite mais pas l'avantage des premières
arrivées, qui rencontrent peu de créanciers et accumulent du capital avant la
densification. Une faible corrélation linéaire globale peut masquer une petite
élite fondatrice. Le contrôle doit porter sur les dates de naissance du top,
leur mortalité, leur durée de résidence et une queue estimée à âge fixé. Un
mélange d'âges peut également produire un corps apparemment exponentiel même si
aucune distribution conditionnelle ne l'est.

\subsection{Capital, revenu et valeur nette}

La distinction conceptuelle est utile mais incomplètement exploitée. Le marché
classe les agents par $w$, la faillite par $\NW$, et la revendication empirique
porte parfois sur le revenu brut. Ces variables peuvent avoir des distributions
différentes. Comme $d_0$ est constant, il translate $\NW$ mais déclenche aussi la
sélection démographique ; il n'est donc pas une simple convention comptable. Le
revenu devrait être rapporté en brut, net, réalisé et variation de richesse. Les
pertes de créances sont une variation de patrimoine, pas un revenu négatif dans
la définition actuelle.

\subsection{Nature du modèle}

M2 est d'abord un modèle de réseau de crédit et de sélection démographique, dont
la variable de sortie est une richesse nette synthétique. Il est moins un modèle
de revenu : aucun salaire, travail, consommation ou processus de revenu autonome
n'est représenté. L'extraction $\sqrt w$ est une production capitalisée. La
distribution observée résulte nécessairement de l'interaction entre bilan,
réseau et survie ; l'interpréter comme loi universelle des revenus dépasserait
la portée du mécanisme.

\section{Résultats de la lecture, limites et prochaines expériences}

Le rapport établit algébriquement la conservation lors d'un prêt, la neutralité
moyenne du choc et la présence d'un rappel autarcique. Il documente
préliminairement que certaines variantes ne meurent pas, que la règle naïve crée
une aristocratie d'âge et qu'E7c peut produire simultanément renouvellement,
population plate et queue compatible avec Pareto. Il n'établit pas le corps
exponentiel ; il le réfute même provisoirement sur le prototype. Il n'établit pas
non plus que la stabilité apparente de la queue provient de la rétroaction SOC.

La première implémentation devra figer une sémantique de référence, tester chaque
invariant, puis reproduire E7c sans chercher à améliorer la forme. Les premières
expériences de réfutation seront : (1) baseline multi-seed et fenêtres séparées ;
(2) corps de $w$, $\NW$ et revenus avec domaines alternatifs ; (3) ablations GBM,
crédit et cascades ; (4) diagnostics anti-cohorte ; (5) variantes d'ordre de
service et de liquidation ; (6) grille $(k,\sigma,\varepsilon_0/w_0)$ incluant
tous les échecs.

\section{Conclusion honnête}

M2 est une proposition scientifique mieux falsifiable que le WIP, non un modèle
validé. La queue de Pareto et le bornage démographique sont plausibles dans une
zone déjà sélectionnée ; le corps exponentiel est absent des résultats rapportés
et l'interprétation SOC reste une hypothèse. La dette d'amorçage, la règle de taux
et la liquidation prorata sont des choix structurants. La suite doit donc viser
une implémentation fidèle et traçable, puis rechercher activement des explications
plus simples --- GBM avec absorption, sélection par âge ou concentration induite
par les règles --- avant d'attribuer la forme à l'auto-organisation critique.

\end{document}
